\section{Conventions and parameters}\label{sec:setup}

Throughout the proof $d\ge3$ is fixed, $q=d-1$, $s=2\sqrt q$, and $n$ is
even.  Constants may depend on $d$ and on the fixed parameters and moment
orders chosen below, but never on $n$.  The assertion that a bound has an
arbitrarily small power loss means that, for each prescribed fixed
$\eta>0$, the bound holds with factor $n^\eta$ once its preceding loss
parameters have been chosen sufficiently small and $n$ is sufficiently
large.  A bound denoted $n^{o(1)}$ has this usual meaning for every fixed
positive exponent.  The two uses will be kept distinct when selecting
parameters.

All vertex-space precision matrices and their inverses act on
$\mathbf1^\perp$.  For coordinate entries and traces we extend such an
operator by zero on the constant vector.  Stub-space operators on a
contrast subspace are extended by zero in the same way.  The distinction
between the identity on a contrast subspace and its zero extension is
understood in all compressed identities.

\subsection{Matrices used during the construction}
The matrix estimates use the inverse of a precision matrix assigned to
the current partial pairing; its normalized form will be denoted by $W$.
The first reference, denoted by $P$, is a finite weighted sum of
nonbacktracking paths, with degrees below the cutoff $L$ fixed below.
Through cleanup the entry error is $E=W-P$. At the separated boundary
we replace this path sum by a smooth truncation $\widehat P_0$ and
continue it by exact inverse updates, obtaining $\widehat P$; the error
is then $E=W-\widehat P$. The precise matrix definitions and update
identities begin in Section~\ref{sec:foundations}. No positivity of
$\widehat P$ is needed.

The symbols below are defined in full where their objects first enter
the proof; the table serves as a reference when moving between stages.
\begin{center}
\begin{tabular}{@{}p{.24\textwidth}p{.70\textwidth}@{}}
\toprule
Notation&Meaning and definition\\
\midrule
$S,l$&Current unresolved stubs and their number; Section~\ref{sec:foundations}.\\
$\Xi,\mathcal D$&Centered incidence and $\mathcal D=\Xi\Xi^T$; Section~\ref{sec:foundations}.\\
$M,W$&Signed precision and its normalized inverse; Section~\ref{sec:foundations}.\\
$P,E$&Truncated path sum and $E=W-P$; Lemma~\ref{lemma:configuration-algebra}.\\
$\widehat P,E$&Continued smooth reference and $E=W-\widehat P$; Section~\ref{sec:boundary-completion}.\\
$\Gamma,D$&Stub compression $\Gamma=\Xi^TW\Xi$ and $D=\Gamma-\Pi_l$; $D$ is distinct from $\mathcal D$.\\
$L_0,L_1,\mathcal T$&Centered traces and $\mathcal T=\operatorname{Tr}(\Xi^TW^2\Xi)$; Section~\ref{subsec:resolvent-consequences}.\\
\bottomrule
\end{tabular}
\end{center}

\subsection{Stopping sizes and spectral parameters}
The unresolved-stub counts decrease as
$dn\to l_1\to l_0\to l_2\to l_3\to0$.
The first two stops delimit early pairing and cleanup; $l_2$ separates
the two boundary phases, and $l_3$ is the final arbitrary matching.
Fix $a=10^{-4}$ and put $e=n^{-2/3+a}$.  The following table records the
principal scales.  The relation $\asymp$ permits fixed positive factors
and integer rounding.
\begin{center}
\begin{tabular}{lll}
\toprule
Quantity&Scale&Use\\
\midrule
$l_0$&$n^{2/3+.15a}$&Separated boundary size\\
$b$&$\lceil n^{1/3+.345a}\rceil$&Cleanup steps\\
$l_1$&$l_0+2b$&End of early pairing\\
$q^{L/2}$&$\sqrt{l_0}\,n^{-.29a}$&Path cutoff\\
$q^{R_0}$&$n^{.04a}$&Boundary separation\\
$q^{g_0}$&$n^{.002a}$&Girth exclusion\\
$\alpha_*$&$n^{-1/3+a/2}$&Closest precision parameter\\
$l_2$&$l_0 n^{-g-r}$&End of boundary phase A\\
$l_3$&$n^\chi$&End of concentrated updates\\
\bottomrule
\end{tabular}
\end{center}
The quantities $l_0,l_1,l_2,l_3$ are even.  A logarithmic radius or cutoff
is rounded to an integer, changing its associated power of $q$ by a fixed
factor.  We keep path degrees $<L$.

For $0<h<1/\sqrt q$ define
\[
 z(h)=h^{-1}+qh,\qquad \alpha(h)=1-qh^2,
 \qquad m_h=\frac{h}{1-h^2}.
\]
Choose the largest primary parameter $h_*$ with
$\alpha(h_*)\asymp n^{-1/3+a/2}$ and
$0<z(h_*)-s<e/20$.  Write $\alpha_* = \alpha(h_*)$.
Successive smaller primary parameters have $1-\sqrt qh$ increasing by a
fixed factor, for instance a factor of two, up to a sufficiently small
fixed constant.  Consequently there are $O(\log n)$ primaries, their
$z-s$ scales cover $[e/20,c]$ up to fixed factors, and successive positive
scales have ratios bounded above and bounded away from one.  Every
parameter remains in a fixed compact subinterval of $(0,1)$.

During the early and cleanup phases, add an auxiliary point immediately
below each primary, at distance
\[
 b'=n^{-.16a}\alpha(h).
\]
At the largest primary also use $h_*-j b_h$, $1\le j\le k_h$, where
\[
 b_h=n^{-.03a}\alpha_*,\qquad
 k_h=\left\lceil\frac{10}{.03a}\right\rceil.
\]
The $j$-spacing is exact.  These auxiliary groups are disjoint for large
$n$.  Both signs $\sigma\in\{1,-1\}$ are used at every tracked parameter.
Only the primary points are retained after $l_0$ unresolved stubs remain.

\subsection{Order of the constant choices}\label{subsec:hierarchy}
The stopping sizes and spectral parameters above specify the stages of
the construction. The remaining constants are introduced with the
estimates they control in Sections~\ref{sec:early-phase}--\ref{sec:boundary-completion}.
The following list records the dependencies among those choices, for
reference when checking that every constant is fixed independently of
$n$. In particular, a coefficient governing growth over a whole phase
must be fixed before the higher moment order used to extract an
individual bound. The list is a dependency guide; the indicated later
arguments define the comparators, row bounds, and exceptional events.
\begin{enumerate}
\item Boundary phase A (Section~\ref{subsec:boundary-phase-a}) uses the
fixed $d,a$ and spectral moment order $k=8$. Its initial row bounds have
constants depending only on $d$; the row comparison constants, including
the finite enlargements at exceptional events, are fixed at this stage.
\item The smooth reference of Section~\ref{subsec:boundary-reference}
has a row-tail exponent $r_P>0$; it is permissible to take $r_P=.001a$.
The exponents $g=r>0$ are chosen sufficiently small relative to $a,r_P$
and the finitely many fixed drift constants in phase A. These constants
depend on $k,d$ and the fixed row bounds, not on the higher scalar moment
orders.
\item In phase A the comparator exponent $\zeta_X$ is chosen as in
\eqref{eq:original-50}, followed by $\theta>0$ sufficiently small
compared with $g$. The row-barrier exponents satisfy
$0<r_v<r_s<r_g<r/2$ with strict gaps. The final matching exponent $\chi$
satisfies $g+r\ll\chi<r_P$ and is larger than these row-barrier exponents.
\item In Section~\ref{sec:cleanup}, the exponent $\rho>0$ is small enough
that $20\rho$ is below the separation and scale margins. For each scalar
concentration account, its normalization exponent is smaller than its
drift and variance margins; its fixed even moment order is then chosen
large enough for extraction over its polynomial-size family.
\item For the early high-moment account of
Lemma~\ref{lem:early-high-moments}, a sufficiently large fixed even $p$
is chosen using the uniform growth coefficient in
\eqref{eq:original-33}. Its diagonal weight and outer-entry threshold
are fixed next, determining the finite near radius $R_c$.
The constants used to choose this threshold depend only on $p,d$, the
finite splitting inequalities, and the uniform row bounds.  Constants
depending on the resulting $R_c$ occur only in terms already accompanied
by a negative power of $n$.
\item Finally, the structural and low-square loss exponents are chosen
small enough for these fixed moment orders and for the finite induction
over near radii. The fixed diagram excess cutoff is large enough for all
the moments that have been used.
\end{enumerate}
There is no dependence of a polynomial horizon-growth coefficient on a
later scalar extraction order.  The estimates below record this fact at
the points where it is essential.  All strict comparisons are unchanged
if scales with rational exponents are replaced by dyadic rational
approximations within fixed factors.  Rational choices of the parameters
$h$ with $O_d(\log n)$ bits, including exact equally spaced auxiliaries,
will be used for the algorithm.
